\documentclass[11pt]{article}

\usepackage{amsmath,amssymb,amsthm}
\usepackage{geometry}
\usepackage{macros}
\geometry{margin=1in}

\title{A Phenomenological Correspondence Between Catuskoti Logic and Binary Logic}
\author{Masaomi Chiba}
\date{2025-12-31}

\begin{document}
\maketitle

\begin{abstract}
This paper presents a phenomenological specification of Catuskoti (fourfold) logic
and binary logic within a unified framework.
The goal is not foundational justification or truth evaluation,
but a descriptive correspondence based on observation, projection, fixed points,
and complexity-based distance measures.
Binary logic is treated as a restricted observational projection
of the same underlying structure described fully by Catuskoti logic.
\end{abstract}

\section{Scope and Assumptions}

The present work adopts the following constraints:

\begin{itemize}
\item No claims of correctness, proof, or logical superiority are made.
\item All descriptions are phenomenological and observational.
\item Truth is treated as a geometric or informational relation, not a boolean value.
\item Dynamics, temporal evolution, and proof theory are explicitly excluded.
\end{itemize}

\section{Shared Underlying Structure}

Both Catuskoti logic and binary logic are assumed to operate on a common underlying
space of propositions or descriptions.
Differences between the two arise solely from:

\begin{itemize}
\item The choice of observational projection.
\item The definition of distance used to identify ``truth.''
\end{itemize}

A distinguished point $\Phi$ is defined in this space and referred to as the
\emph{middle-way fixed point}.

\section{Correspondence Table}

\begin{center}
\begin{tabular}{|l|l|l|}
\hline
Observation Aspect & Catuskoti Logic & Binary Logic \\
\hline
Underlying space & Proposition space & Same space \\
Observation projection & Fourfold projection & Two-valued projection \\
Accessible phases & Is, Not-Is, Both, Neither & Is / Not-Is only \\
Fixed point & Middle-way fixed point $\Phi$ & Same fixed point $\Phi$ \\
Fixed point behavior & Stable & Repulsive \\
Truth definition & Minimum Kolmogorov distance to $\Phi$ & Maximum distance from infinity \\
Truth behavior & Convergence to $\Phi$ & Divergence toward extremes \\
Undecidability & Internal phase & Explicit indeterminacy \\
Observation outside projection & None & Intermediate phases \\
\hline
\end{tabular}
\end{center}

\section{Incompleteness as an Observational Effect}

Incompleteness phenomena are interpreted as consequences of projection choice,
not as failures of the underlying structure.

\subsection{Godel-Type Phenomena}

In binary logic, certain propositions become undecidable due to the restriction
to two-valued observation.
In Catuskoti logic, the same propositions occupy intermediate observational phases
and remain representable.

\subsection{Chaitin-Type Phenomena}

When truth is defined via extremal distance criteria,
complexity bounds appear as external limitations.
Under minimum-distance evaluation, these bounds are treated as observable phases
rather than obstructions.

\section{Interpretive Summary}

\begin{itemize}
\item Catuskoti logic provides a complete observational decomposition of the
proposition space relative to the fixed point $\Phi$.
\item Binary logic corresponds to a coarse projection of the same space.
\item Both systems share the same fixed point, but assign it opposite stability roles.
\item Incompleteness arises from projection, not from the space itself.
\end{itemize}

\section{Conclusion}

This specification demonstrates that Catuskoti logic and binary logic
may be treated as isomorphic descriptions of a single underlying structure,
differing only in observational resolution and distance-based truth criteria.
No dynamical, proof-theoretic, or metaphysical claims are required for this correspondence.

\end{document}

